%/*% pdflatex -file-line-error -interaction nonstopmode %

\documentclass{article}
\usepackage{graphicx}
\usepackage{amsmath}
\usepackage{amsfonts}
\usepackage{bookmark}
\usepackage{exercises}
\usepackage{amsthm}
\usepackage{yhmath}
\newcommand*{\oldneg}{\mathord{\sim}}

\theoremstyle{definition}
\newtheorem{definition}{Definition}

\setlength{\parindent}{.5in}
\setlength{\parskip}{12pt}

\title{Misc Exercises}
\begin{document}
\maketitle

\section*{Proof of the Similarity of Parabolas}
We consider the vertex form of any quadratic function, $y=a(x-h)^2+k$. Changing $x$ and $y$ does not change the steepnees of the curve of the parabola, but only the 
vertex (a fact that we decided should wait to be proven until we at least learn a bit about complex numbers), and hence, the actual curvature is given by $y=ax^2$. Now, 
consider any other $a$, let it be $m$; hence $m=al$ for some $l$. Well, we multiply both $x$ and $y$ in the above equation by $l$, thus recieving:\newline
$yl=a(xl)^2\newline\rightarrow yl=ax^2l^2\newline\rightarrow y=alx^2$. Thus, we have effectively gone from every quadratic equation in vertex form by scaling the 
$x$ and $y$ coordinates, which proves that all quadratic equations are similar.
\end{document}
